\chapter{Independent Sets Cliques}
%%%%%%%%%%%%%%%%%%%%%%%%%%%%%%%%%%%%%%%%%%%%%%%%%%%%%%%%%%%%%%%%%%%%%%%%%%%%%%%
\section{Independent Sets Cliques}

% Generated by scripts/blueprint_restate.py from TCSlib.GraphTheory.Core.IndependentSetsCliques.
% Each body is the STATEMENT only, taken from the declaration's Lean docstring:
% the one-line summary plus the verbatim book statement.  Proof, proof plan,
% significance commentary and formalisation notes are excluded by construction.

\begin{definition}[A covering $K$: every edge has an end in $K$]
\lean{SimpleGraph.IsCovering}
\label{SimpleGraph.IsCovering}
\leanok
A \textbf{covering} \texttt{K}: every edge has an end in \texttt{K}.

Recall that a subset $K$ of $V$ such that every edge of $G$ has at least one end
in $K$ is called a covering of $G$.
\end{definition}

\begin{theorem}[Independent sets are complements of vertex coverings]
\lean{SimpleGraph.isIndepSet_iff_isCovering_compl}
\label{SimpleGraph.isIndepSet_iff_isCovering_compl}
\uses{SimpleGraph.IsCovering}
Let $G$ be a graph with vertex set $V$, and let $S \subseteq V$. Then $S$ is an
independent set of $G$ — no edge of $G$ has both of its ends in $S$ — if and only if the
complement $V \setminus S$ is a covering of $G$, that is, every edge of $G$ has at least
one end in $V \setminus S$.
\end{theorem}

\begin{definition}[The covering number $\beta(G)$]
\lean{SimpleGraph.coveringNumber}
\label{SimpleGraph.coveringNumber}
\leanok
\uses{SimpleGraph.IsCovering}
The covering number \texttt{$\beta$(G)}. No \texttt{Nat.sInf $\emptyset$ = 0} trap:
\texttt{Set.univ} always covers, so the set is nonempty.

The number of vertices in a maximum independent set of $G$ is called the
\emph{independence number} of $G$ and is denoted by $\alpha(G)$; similarly, the
number of vertices in a minimum covering of $G$ is the \emph{covering number} of $G$
and is denoted by $\beta(G)$.
\end{definition}

\begin{theorem}[Gallai's identity relating independence and covering numbers]
\lean{SimpleGraph.indepNum_add_coveringNumber}
\label{SimpleGraph.indepNum_add_coveringNumber}
\uses{SimpleGraph.coveringNumber}
Let $G$ be a simple graph on a finite vertex set $V$, and write $n = \abs{V}$ for its
number of vertices. Denote by $\alpha(G)$ the independence number of $G$, the largest
size of a set of vertices no two of which are adjacent, and by $\beta(G)$ the covering
number of $G$, the smallest size of a set of vertices meeting every edge. Then \[
\alpha(G) + \beta(G) = n. \]
\end{theorem}

\begin{definition}[The matching number $\alpha'(G)$]
\lean{SimpleGraph.matchingNumber}
\label{SimpleGraph.matchingNumber}
\leanok
The matching number \texttt{$\alpha$$'$(G)}.

The edge analogue of an independent set is a set of links no two of which are
adjacent, that is, a matching. [...] We denote the number of edges in a maximum
matching of $G$ by $\alpha'(G)$ [...] the \emph{edge independence number}.
\end{definition}

\begin{definition}[An edge covering $L$: every vertex is an end of some edge of $L$]
\lean{SimpleGraph.IsEdgeCovering}
\label{SimpleGraph.IsEdgeCovering}
\leanok
An \textbf{edge covering} \texttt{L}: every vertex is an end of some edge of \texttt{L}.

The edge analogue of a covering is called an edge covering. An \emph{edge covering} of
$G$ is a subset $L$ of $E$ such that each vertex of $G$ is an end of some edge in $L$.
Note that edge coverings do not always exist; a graph $G$ has an edge covering if and
only if $\delta > 0$.
\end{definition}

\begin{definition}[The edge covering number $\beta'(G)$]
\lean{SimpleGraph.edgeCoveringNumber}
\label{SimpleGraph.edgeCoveringNumber}
\leanok
\uses{SimpleGraph.IsEdgeCovering}
The edge covering number \texttt{$\beta$$'$(G)}. \texttt{Nat.sInf $\emptyset$ = 0} IS a
live trap when \texttt{$\delta$ = 0} --- hence the \texttt{$\delta$ > 0} hypothesis in
Thm 7.2.

[We denote] the number of edges in a minimum edge covering of $G$ by
$\beta'(G)$; the numbers $\alpha'(G)$ and $\beta'(G)$ are the \emph{edge
independence number} and \emph{edge covering number} of $G$, respectively.
\end{definition}

\begin{theorem}[Gallai's identity for matching and edge covering numbers]
\lean{SimpleGraph.matchingNumber_add_edgeCoveringNumber}
\label{SimpleGraph.matchingNumber_add_edgeCoveringNumber}
\uses{SimpleGraph.edgeCoveringNumber, SimpleGraph.matchingNumber}
Let $G$ be a finite simple graph on $n$ vertices with minimum degree $\delta(G) > 0$ (so
that $G$ has no isolated vertices, and edge coverings exist). Then the matching number
and the edge covering number of $G$ satisfy
\[
\alpha'(G) + \beta'(G) = n,
\]
where $\alpha'(G)$ is the largest number of edges in a matching of $G$ and $\beta'(G)$
is the smallest number of edges in an edge covering of $G$.

This is the theorem of Gallai (1959).
\end{theorem}

\begin{theorem}[Independence number equals edge covering number in bipartite graphs]
\lean{SimpleGraph.indepNum_eq_edgeCoveringNumber}
\label{SimpleGraph.indepNum_eq_edgeCoveringNumber}
\uses{SimpleGraph.edgeCoveringNumber}
Let $G$ be a finite bipartite graph with vertex classes $X$ and $Y$, and suppose that
$G$ has minimum degree $\delta(G) > 0$, so that every vertex is an endpoint of at least
one edge. Then the independence number $\alpha(G)$—the greatest number of pairwise
non-adjacent vertices—equals the edge covering number $\beta'(G)$, the least number of
edges in an edge covering of $G$ (a set of edges meeting every vertex).
\end{theorem}

\begin{theorem}[Bipartiteness via independence numbers of subgraphs]
\lean{SimpleGraph.isBipartite_iff_forall_subgraph_two_mul_indepNum_le}
\label{SimpleGraph.isBipartite_iff_forall_subgraph_two_mul_indepNum_le}
Let $G$ be a finite simple graph. Then $G$ is bipartite if and only if every subgraph
$H$ of $G$ satisfies
\[
\abs{V(H)} \le 2\,\alpha(H),
\]
where $V(H)$ is the vertex set of $H$ and $\alpha(H)$ denotes the independence number of
$H$, that is, the size of a largest set of pairwise non-adjacent vertices of $H$.
\end{theorem}

\begin{theorem}[Bipartiteness via independence and edge covering numbers]
\lean{SimpleGraph.isBipartite_iff_forall_subgraph_indepNum_eq_edgeCoveringNumber}
\label{SimpleGraph.isBipartite_iff_forall_subgraph_indepNum_eq_edgeCoveringNumber}
\uses{SimpleGraph.edgeCoveringNumber}
A finite simple graph $G$ is bipartite if and only if, for every subgraph $H$ of $G$
whose minimum degree $\delta(H)$ is positive, the independence number $\alpha(H)$ — the
greatest number of pairwise nonadjacent vertices of $H$ — equals the edge covering
number $\beta'(H)$.
\end{theorem}

\begin{definition}[$\alpha$-critical: deleting any edge raises the independence number]
\lean{SimpleGraph.IsAlphaCritical}
\label{SimpleGraph.IsAlphaCritical}
\leanok
$\alpha$-critical: deleting any edge raises the independence number.

7.1.2   A graph is $\alpha$-\emph{critical} if $\alpha(G - e) > \alpha(G)$ for
all $e \in E$.
\end{definition}

\begin{theorem}[Alpha-critical connected graphs have no cut vertex]
\lean{SimpleGraph.no_cutVertex_of_isAlphaCritical}
\label{SimpleGraph.no_cutVertex_of_isAlphaCritical}
\uses{SimpleGraph.IsAlphaCritical, SimpleGraph.IsVertexCut}
Let $G$ be a finite simple graph, and write $\alpha(G)$ for its independence number (the
size of a largest set of pairwise non-adjacent vertices). Suppose $G$ is connected and
$\alpha$-critical, meaning that $\alpha(G - e) > \alpha(G)$ for every edge $e$ of $G$.
Then $G$ has no cut vertex: for every vertex $v$, the singleton $\{v\}$ is not a vertex
cut, that is, it is not the case that $\{v\}$ is a proper subset of the vertex set whose
deletion leaves the graph induced on the remaining vertices disconnected.
\end{theorem}

\begin{definition}[$\beta$-critical: deleting any edge lowers the covering number]
\lean{SimpleGraph.IsBetaCritical}
\label{SimpleGraph.IsBetaCritical}
\leanok
\uses{SimpleGraph.coveringNumber}
$\beta$-critical: deleting any edge lowers the covering number.

7.1.3   A graph $G$ is $\beta$-\emph{critical} if $\beta(G - e) < \beta(G)$ for
all $e \in E$.
\end{definition}

\begin{theorem}[Beta-critical connected graphs have no cut vertex]
\lean{SimpleGraph.no_cutVertex_of_isBetaCritical}
\label{SimpleGraph.no_cutVertex_of_isBetaCritical}
\uses{SimpleGraph.IsBetaCritical, SimpleGraph.IsVertexCut}
Let $G$ be a finite connected simple graph that is $\beta$-critical, meaning that
deleting any single edge strictly decreases its covering number $\beta(G)$ (the least
number of vertices meeting every edge). Then no single vertex is a cut vertex: for every
vertex $v$, the singleton $\{v\}$ is not a vertex cut, that is, removing $v$ never
leaves the graph disconnected.
\end{theorem}

\begin{theorem}[Twice the covering number is at most one more than the edge count]
\lean{SimpleGraph.two_mul_coveringNumber_le}
\label{SimpleGraph.two_mul_coveringNumber_le}
\uses{SimpleGraph.coveringNumber}
Let $G$ be a connected finite simple graph with $\varepsilon$ edges, and let $\beta(G)$
denote its covering number, the least number of vertices in a covering of $G$. Then \[
2\,\beta(G) \le \varepsilon + 1. \]
\end{theorem}

\begin{definition}[\texttt{IsRamseyBound n k l}: every graph on $n$ vertices has a $k$-clique or an $l$-indep set.\dots]
\lean{SimpleGraph.IsRamseyBound}
\label{SimpleGraph.IsRamseyBound}
\leanok
\texttt{IsRamseyBound n k l}: every graph on \texttt{n} vertices has a \texttt{k}-clique or an \texttt{l}-indep set.
\end{definition}

\begin{definition}[The Ramsey number $r(k, l)$]
\lean{SimpleGraph.ramseyNumber}
\label{SimpleGraph.ramseyNumber}
\leanok
\uses{SimpleGraph.IsRamseyBound}
The Ramsey number \texttt{r(k, l)}. \texttt{Nat.sInf $\emptyset$ = 0} is a live trap
until \texttt{exists\_isRamseyBound} is proved.

The numbers $r(k, l)$ are known as the \emph{Ramsey numbers}.

with the defining property quoted on \texttt{IsRamseyBound} above.
\end{definition}

\begin{theorem}[Ramsey bounds are upward closed in the number of vertices]
\lean{SimpleGraph.isRamseyBound_mono}
\label{SimpleGraph.isRamseyBound_mono}
\uses{SimpleGraph.IsRamseyBound}
Fix natural numbers $k$ and $l$, and say that $n$ is a Ramsey bound for $(k,l)$ if every
simple graph on $n$ vertices contains either a clique on $k$ vertices or an independent
set on $l$ vertices. If $n \le m$ and $n$ is a Ramsey bound for $(k,l)$, then $m$ is
also a Ramsey bound for $(k,l)$.
\end{theorem}

\begin{theorem}[Existence of Ramsey numbers]
\lean{SimpleGraph.exists_isRamseyBound}
\label{SimpleGraph.exists_isRamseyBound}
\uses{SimpleGraph.IsRamseyBound}
For every pair of natural numbers $k$ and $l$, there exists a natural number $n$ with
the following property: every simple graph on $n$ vertices contains either a clique of
size $k$ or an independent set of size $l$.
\end{theorem}

\begin{theorem}[The Ramsey number attains its bound]
\lean{SimpleGraph.isRamseyBound_ramseyNumber}
\label{SimpleGraph.isRamseyBound_ramseyNumber}
\uses{SimpleGraph.IsRamseyBound, SimpleGraph.ramseyNumber}
For natural numbers $k$ and $l$, let $r(k,l)$ denote the Ramsey number, defined as the
least $n$ (equivalently, the infimum over all such $n$) for which every simple graph on
$n$ vertices contains either a clique of size $k$ or an independent set of size $l$.
Then $r(k,l)$ itself has this property: every simple graph on $r(k,l)$ vertices contains
a clique of size $k$ or an independent set of size $l$.
\end{theorem}

\begin{theorem}[The Ramsey number $r(1,l)$]
\lean{SimpleGraph.ramseyNumber_one_left}
\label{SimpleGraph.ramseyNumber_one_left}
\uses{SimpleGraph.ramseyNumber}
Let the Ramsey number $r(k,l)$ denote the least natural number $n$ such that every
simple graph on $n$ vertices contains either a clique on $k$ vertices or an independent
set on $l$ vertices. Then for every integer $l \ge 1$ one has $r(1,l) = 1$.
\end{theorem}

\begin{theorem}[The Ramsey number $r(k,1)$ equals one]
\lean{SimpleGraph.ramseyNumber_one_right}
\label{SimpleGraph.ramseyNumber_one_right}
\uses{SimpleGraph.ramseyNumber}
For every integer $k \ge 1$, the Ramsey number satisfies $r(k, 1) = 1$; that is, $1$ is
the least number of vertices $n$ for which every simple graph on $n$ vertices contains
either a clique of size $k$ or an independent set of size $1$.
\end{theorem}

\begin{theorem}[The two-color Ramsey number with clique size two]
\lean{SimpleGraph.ramseyNumber_two_left}
\label{SimpleGraph.ramseyNumber_two_left}
\uses{SimpleGraph.ramseyNumber}
For every integer $l \ge 1$, the Ramsey number satisfies $r(2, l) = l$. Equivalently,
$l$ is the least number of vertices $n$ such that every simple graph on $n$ vertices
contains either an edge (a clique on $2$ vertices) or an independent set of size $l$.
\end{theorem}

\begin{theorem}[The Ramsey number $r(k,2)$]
\lean{SimpleGraph.ramseyNumber_two_right}
\label{SimpleGraph.ramseyNumber_two_right}
\uses{SimpleGraph.ramseyNumber}
Write $r(k,l)$ for the Ramsey number: the least $n$ for which every simple graph on $n$
vertices contains either a clique of size $k$ or an independent set of size $l$. Then
for every integer $k \ge 1$,
\[
r(k, 2) = k.
\]
\end{theorem}

\begin{theorem}[The Ramsey recursion inequality]
\lean{SimpleGraph.ramsey_recursion}
\label{SimpleGraph.ramsey_recursion}
\uses{SimpleGraph.ramseyNumber}
For integers $k \ge 2$ and $l \ge 2$, the Ramsey numbers satisfy
\[
r(k, l) \le r(k, l-1) + r(k-1, l),
\]
where $r(k, l)$ denotes the least number of vertices $n$ such that every simple graph on
$n$ vertices contains either a clique on $k$ vertices or an independent set on $l$
vertices.
\end{theorem}

\begin{theorem}[Strict Ramsey recurrence for even summands]
\lean{SimpleGraph.ramsey_recursion_strict}
\label{SimpleGraph.ramsey_recursion_strict}
\uses{SimpleGraph.ramseyNumber}
For natural numbers $k, l \ge 2$, let $r(k, l)$ denote the Ramsey number: the least $n$
such that every simple graph on $n$ vertices contains either a clique on $k$ vertices or
an independent set on $l$ vertices. If both $r(k, l-1)$ and $r(k-1, l)$ are even, then
\[
r(k, l) < r(k, l-1) + r(k-1, l).
\]
\end{theorem}

\begin{theorem}[The five-cycle has no triangle]
\lean{SimpleGraph.cycleGraph_five_cliqueFree_three}
\label{SimpleGraph.cycleGraph_five_cliqueFree_three}
The cycle graph $C_5$ on five vertices contains no clique of size $3$; that is, no three
of its vertices are pairwise adjacent.
\end{theorem}

\begin{theorem}[The 5-cycle has no independent set of three vertices]
\lean{SimpleGraph.cycleGraph_five_indepSetFree_three}
\label{SimpleGraph.cycleGraph_five_indepSetFree_three}
Let $C_5$ denote the cycle graph on five vertices, in which the vertices are arranged in
a ring and each vertex is adjacent exactly to its two neighbours. Then $C_5$ contains no
independent set of three vertices; that is, no three of its vertices are pairwise
non-adjacent.
\end{theorem}

\begin{theorem}[A lower bound for the Ramsey number $r(3,3)$]
\lean{SimpleGraph.six_le_ramseyNumber_three_three}
\label{SimpleGraph.six_le_ramseyNumber_three_three}
\uses{SimpleGraph.ramseyNumber}
Let $r(3,3)$ denote the Ramsey number for the pair $(3,3)$: the least natural number $n$
such that every simple graph on $n$ vertices contains either three mutually adjacent
vertices (a triangle) or three mutually non-adjacent vertices (an independent set of
size $3$). Then
\[
r(3,3) \ge 6.
\]
\end{theorem}

\begin{definition}[The (3,5)-Ramsey graph: \texttt{ZMod 13}, adjacent iff the difference is a cubic residue]
\lean{SimpleGraph.cubicResidueGraph13}
\label{SimpleGraph.cubicResidueGraph13}
The (3,5)-Ramsey graph: \texttt{ZMod 13}, adjacent iff the difference is a cubic
residue. \texttt{\{1,5,8,12\}} is closed under negation, so \texttt{symm} is genuine.

A $(k, l)$-\emph{Ramsey graph} is a graph on $r(k, l) - 1$ vertices that contains
neither a clique of $k$ vertices nor an independent set of $l$ vertices. [...]
Ramsey graphs often seem to possess interesting structures. [...] We get the
$(3, 5)$-Ramsey graph by regarding the thirteen vertices as elements of the
field of integers modulo 13, and joining two vertices by an edge if their
difference is a cubic residue of 13 (either 1, 5, 8 or 12).
\end{definition}

\begin{theorem}[Triangle-freeness of the cubic-residue graph mod 13]
\lean{SimpleGraph.cubicResidueGraph13_cliqueFree_three}
\label{SimpleGraph.cubicResidueGraph13_cliqueFree_three}
\uses{SimpleGraph.cubicResidueGraph13}
Consider the graph on the thirteen vertices $\bbz/13\bbz$ in which two distinct vertices
are joined by an edge exactly when their difference is a nonzero cubic residue modulo
$13$, that is, an element of $\{1,5,8,12\}$. This graph contains no clique on three
vertices: there is no triple of pairwise adjacent vertices, so the graph is
triangle-free.
\end{theorem}

\begin{theorem}[No independent set of five in the cubic-residue graph mod 13]
\lean{SimpleGraph.cubicResidueGraph13_indepSetFree_five}
\label{SimpleGraph.cubicResidueGraph13_indepSetFree_five}
\uses{SimpleGraph.cubicResidueGraph13}
Let $G$ be the graph on the thirteen vertices $\bbz/13\bbz$ in which two distinct
vertices are joined by an edge exactly when their difference is a cubic residue modulo
$13$, that is, an element of $\{1,5,8,12\}$. Then $G$ contains no independent set of
five vertices: among any five vertices, some two are adjacent.
\end{theorem}

\begin{theorem}[Lower bound on the Ramsey number $r(3,5)$]
\lean{SimpleGraph.fourteen_le_ramseyNumber_three_five}
\label{SimpleGraph.fourteen_le_ramseyNumber_three_five}
\uses{SimpleGraph.ramseyNumber}
Let $r(k,l)$ denote the Ramsey number, the least $n$ such that every simple graph on $n$
vertices contains either a clique of size $k$ or an independent set of size $l$. Then
\[r(3,5) \ge 14.\]
\end{theorem}

\begin{definition}[The (4,4)-Ramsey graph = the Paley graph of order 17]
\lean{SimpleGraph.paleyGraph17}
\label{SimpleGraph.paleyGraph17}
The (4,4)-Ramsey graph = the Paley graph of order 17.

the $(4, 4)$-Ramsey graph is obtained by regarding the vertices as elements of the field
of integers modulo 17, and joining two vertices if their difference is a quadratic
residue of 17 (either 1, 2, 4, 8, 9, 13, 15 or 16). It has been conjectured that the
$(k, k)$-Ramsey graphs are always self-complementary (that is, isomorphic to their
complements); this is true for $k = 2$, 3 and 4.
\end{definition}

\begin{theorem}[Absence of a 4-clique in the Paley graph of order 17]
\lean{SimpleGraph.paleyGraph17_cliqueFree_four}
\label{SimpleGraph.paleyGraph17_cliqueFree_four}
\uses{SimpleGraph.paleyGraph17}
Let $P_{17}$ denote the Paley graph of order $17$: its vertices are the elements of the
field $\bbz/17\bbz$, and two vertices are joined by an edge whenever their difference is
a nonzero quadratic residue modulo $17$ (that is, one of $1, 2, 4, 8, 9, 13, 15, 16$).
Then $P_{17}$ contains no clique on four vertices; equivalently, there is no set of four
pairwise adjacent vertices.
\end{theorem}

\begin{theorem}[The Paley graph on 17 vertices has no independent set of size four]
\lean{SimpleGraph.paleyGraph17_indepSetFree_four}
\label{SimpleGraph.paleyGraph17_indepSetFree_four}
\uses{SimpleGraph.paleyGraph17}
Let $G$ be the Paley graph of order $17$: its vertices are the elements of the field
$\bbz/17\bbz$, and two vertices are joined by an edge exactly when their difference is a
nonzero quadratic residue modulo $17$ (that is, one of $1, 2, 4, 8, 9, 13, 15, 16$).
Then $G$ contains no independent set of size $4$; equivalently, among any four vertices
of $G$ some two are adjacent.
\end{theorem}

\begin{theorem}[Lower bound on the Ramsey number $r(4,4)$]
\lean{SimpleGraph.eighteen_le_ramseyNumber_four_four}
\label{SimpleGraph.eighteen_le_ramseyNumber_four_four}
\uses{SimpleGraph.ramseyNumber}
Write $r(4,4)$ for the least $n$ such that every simple graph on $n$ vertices contains
either a clique on $4$ vertices or an independent set of $4$ vertices. Then $r(4,4) \ge
18$.
\end{theorem}

\begin{theorem}[Erdős–Szekeres bound on Ramsey numbers]
\lean{SimpleGraph.ramseyNumber_le_choose}
\label{SimpleGraph.ramseyNumber_le_choose}
\uses{SimpleGraph.ramseyNumber}
For integers $k, l \ge 1$, the Ramsey number $r(k, l)$ — the least $n$ such that every
simple graph on $n$ vertices contains either a clique on $k$ vertices or an independent
set on $l$ vertices — satisfies
\[
r(k, l) \le \binom{k + l - 2}{k - 1}.
\]
\end{theorem}

\begin{theorem}[Number of simple graphs on a finite vertex set]
\lean{SimpleGraph.card_simpleGraph_fin}
\label{SimpleGraph.card_simpleGraph_fin}
For every natural number $n$, the number of simple graphs on the vertex set $\{1, 2,
\ldots, n\}$ equals $2^{\binom{n}{2}}$, since a simple graph is determined by choosing
which of the $\binom{n}{2}$ possible unordered pairs of distinct vertices form edges.
\end{theorem}

\begin{theorem}[Erdős lower bound for the diagonal Ramsey number]
\lean{SimpleGraph.ramsey_self_lower_bound}
\label{SimpleGraph.ramsey_self_lower_bound}
\uses{SimpleGraph.ramseyNumber}
For an integer $k \ge 1$, let $r(k,k)$ denote the Ramsey number: the least $n$ such that
every simple graph on $n$ vertices contains either a clique of size $k$ or an
independent set of size $k$. Then
\[
2^{k/2} \le r(k,k).
\]
\end{theorem}

\begin{theorem}[Exponential lower bound for the Ramsey numbers]
\lean{SimpleGraph.ramseyNumber_ge_of_min}
\label{SimpleGraph.ramseyNumber_ge_of_min}
\uses{SimpleGraph.ramseyNumber}
Let $k$ and $l$ be positive integers, and let $r(k,l)$ denote the associated Ramsey
number. Then
\[
2^{\min(k,l)/2} \le r(k,l).
\]
\end{theorem}

\begin{definition}[An $m$-edge-colouring of $K_{n}$: a total function on unordered pairs]
\lean{SimpleGraph.EdgeColouring}
\label{SimpleGraph.EdgeColouring}
\leanok
An \texttt{m}-edge-colouring of \texttt{K$_{n}$}: a total function on unordered pairs.
NOT \texttt{SimpleGraph.Coloring}, which colours vertices.
\end{definition}

\begin{definition}[\texttt{IsRamseyBoundMulti n c}: every $m$-edge-colouring of $K_{n}$ has, for some $i$, a set of $c i$\dots]
\lean{SimpleGraph.IsRamseyBoundMulti}
\label{SimpleGraph.IsRamseyBoundMulti}
\leanok
\uses{SimpleGraph.EdgeColouring}
\texttt{IsRamseyBoundMulti n c}: every \texttt{m}-edge-colouring of \texttt{K$_{n}$} has, for some \texttt{i}, a set of \texttt{c i}
vertices whose internal edges are all colour \texttt{i}.

We define $r(k_1, k_2, \ldots, k_m)$ to be the smallest integer $n$ such that
every $m$-edge colouring $(E_1, E_2, \ldots, E_m)$ of $K_n$ contains, for some
$i$, a complete subgraph on $k_i$ vertices, all of whose edges are in colour
$i$.
\end{definition}

\begin{definition}[The multicolour Ramsey number $r(k_{1},\dots,k_m)$]
\lean{SimpleGraph.ramseyNumberMulti}
\label{SimpleGraph.ramseyNumberMulti}
\leanok
\uses{SimpleGraph.IsRamseyBoundMulti}
The multicolour Ramsey number \texttt{r(k$_{1}$,\dots{},k\_m)}.

The least \texttt{n} for which every \texttt{m}-edge colouring of \texttt{K$_{n}$}
yields a monochromatic
complete subgraph on \texttt{k$_{i}$} vertices in some colour \texttt{i} --- see the
verbatim text on
\texttt{IsRamseyBoundMulti} above.
\end{definition}

\begin{theorem}[A recursive upper bound for multicolour Ramsey numbers]
\lean{SimpleGraph.ramseyMulti_recursion}
\label{SimpleGraph.ramseyMulti_recursion}
\uses{SimpleGraph.ramseyNumberMulti}
Fix an integer $m \ge 1$ of colours and integers $k_1, k_2, \dots, k_m$ with $k_i \ge 2$
for every $i$, and let $r(k_1, \dots, k_m)$ denote the multicolour Ramsey number: the
least $n$ such that every colouring of the edges of the complete graph $K_n$ with $m$
colours contains, for some colour $i$, a complete subgraph on $k_i$ vertices all of
whose edges have colour $i$. Then
\[
r(k_1, \dots, k_m) + m \;\le\; \sum_{i=1}^{m} r(k_1, \dots, k_{i-1}, \, k_i - 1, \,
k_{i+1}, \dots, k_m) \;+\; 2,
\]
where the $i$th summand is obtained from $r(k_1, \dots, k_m)$ by decreasing the argument
$k_i$ by one.
\end{theorem}

\begin{theorem}[Multinomial bound for multicolour Ramsey numbers]
\lean{SimpleGraph.ramseyMulti_le_multinomial}
\label{SimpleGraph.ramseyMulti_le_multinomial}
\uses{SimpleGraph.ramseyNumberMulti}
For any natural numbers $k_1, \ldots, k_m$, the multicolour Ramsey number satisfies
\[
r(k_1 + 1, \ldots, k_m + 1) \;\le\; \frac{(k_1 + k_2 + \cdots + k_m)!}{k_1!\, k_2!
\cdots k_m!},
\]
the right-hand side being the multinomial coefficient $\binom{k_1 + \cdots + k_m}{k_1,
\ldots, k_m}$.
\end{theorem}

\begin{theorem}[Symmetry of the Ramsey numbers]
\lean{SimpleGraph.ramseyNumber_comm}
\label{SimpleGraph.ramseyNumber_comm}
\uses{SimpleGraph.ramseyNumber}
For natural numbers $k$ and $l$, let $r(k,l)$ denote the Ramsey number: the least $n$
such that every simple graph on $n$ vertices contains either a clique on $k$ vertices or
an independent set on $l$ vertices. Then $r(k,l) = r(l,k)$.
\end{theorem}

\begin{definition}[$r_{n} = r(3, \dots, 3)$ with $n$ colours]
\lean{SimpleGraph.ramseyTriangle}
\label{SimpleGraph.ramseyTriangle}
\leanok
\uses{SimpleGraph.ramseyNumberMulti}
\texttt{r$_{n}$ = r(3, \dots{}, 3)} with \texttt{n} colours.

\textbf{7.2.3} Let $r_n$ denote the Ramsey number $r(k_1, k_2, \ldots, k_n)$ with
$k_i = 3$ for all $i$.
\end{definition}

\begin{theorem}[A recursive bound for triangle Ramsey numbers]
\lean{SimpleGraph.ramseyTriangle_recursion}
\label{SimpleGraph.ramseyTriangle_recursion}
\uses{SimpleGraph.ramseyTriangle}
For an integer $n \ge 1$, write $r_n$ for the $n$-colour triangle Ramsey number: the
least $N$ such that every colouring of the edges of the complete graph $K_N$ with $n$
colours contains a triangle whose three edges all receive the same colour. Then
\[
r_n + n \le n\,r_{n-1} + 2 .
\]
\end{theorem}

\begin{theorem}[A factorial bound for the multicolour triangle Ramsey number]
\lean{SimpleGraph.ramseyTriangle_le_factorial_exp}
\label{SimpleGraph.ramseyTriangle_le_factorial_exp}
\uses{SimpleGraph.ramseyTriangle}
For every integer $n \ge 2$, the $n$-colour triangle Ramsey number $r_n = r(3, 3, \dots,
3)$ — the least $N$ for which every colouring of the edges of the complete graph on $N$
vertices with $n$ colours yields a monochromatic triangle in some colour — satisfies
\[
  r_n \le \lfloor n!\, e \rfloor + 1,
\]
where $e = \exp(1)$ and $\lfloor\,\cdot\,\rfloor$ is the floor function.
\end{theorem}

\begin{theorem}[Three-colour triangle Ramsey number is at most 17]
\lean{SimpleGraph.ramseyTriangle_three_le}
\label{SimpleGraph.ramseyTriangle_three_le}
\uses{SimpleGraph.ramseyTriangle}
Let $r_3 = r(3,3,3)$ denote the three-colour Ramsey number for triangles, that is, the
least $n$ such that every colouring of the edges of the complete graph $K_n$ with three
colours contains a monochromatic triangle. Then $r_3 \le 17$; equivalently, every
$3$-edge-colouring of $K_{17}$ admits three vertices whose three connecting edges all
receive the same colour.
\end{theorem}

\begin{definition}[The composition (lexicographic product) $G(H)$]
\lean{SimpleGraph.composition}
\label{SimpleGraph.composition}
The composition (lexicographic product) \texttt{G[H]}.

7.2.4 The \emph{composition} of simple graphs $G$ and $H$ is the simple graph $G[H]$
with vertex set $V(G) \times V(H)$, in which $(u, v)$ is adjacent to
$(u', v')$ if and only if either $uu' \in E(G)$ or $u = u'$ and
$vv' \in E(H)$.
\end{definition}

\begin{theorem}[Independence number of a lexicographic product]
\lean{SimpleGraph.indepNum_composition_le}
\label{SimpleGraph.indepNum_composition_le}
\uses{SimpleGraph.composition}
Let $G$ and $H$ be simple graphs on finite vertex sets, and let $G[H]$ denote their
composition (lexicographic product): the simple graph with vertex set $V(G) \times V(H)$
in which $(u,v)$ is adjacent to $(u',v')$ exactly when $uu' \in E(G)$, or else $u = u'$
and $vv' \in E(H)$. Writing $\alpha(\cdot)$ for the independence number, the
independence number of the composition satisfies \[ \alpha(G[H]) \le
\alpha(G)\,\alpha(H). \]
\end{theorem}

\begin{theorem}[Product lower bound for diagonal Ramsey numbers]
\lean{SimpleGraph.ramsey_product_lower}
\label{SimpleGraph.ramsey_product_lower}
\uses{SimpleGraph.ramseyNumber}
For positive integers $k$ and $l$, write $r(a,b)$ for the Ramsey number: the least $n$
such that every simple graph on $n$ vertices contains either a clique on $a$ vertices or
an independent set of $b$ vertices. Then, provided $k \ge 1$ and $l \ge 1$,
\[
\bigl(r(k+1,\,k+1) - 1\bigr)\bigl(r(l+1,\,l+1) - 1\bigr) + 1 \;\le\; r(kl+1,\,kl+1).
\]
\end{theorem}

\begin{theorem}[Abbott's lower bound for diagonal Ramsey numbers]
\lean{SimpleGraph.abbott_lower_bound}
\label{SimpleGraph.abbott_lower_bound}
\uses{SimpleGraph.ramseyNumber}
For every natural number $n \ge 0$, the diagonal Ramsey number satisfies
\[
  r\bigl(2^n + 1,\; 2^n + 1\bigr) \ge 5^n + 1,
\]
where $r(k, l)$ denotes the least number of vertices $N$ such that every graph on $N$
vertices contains either a clique on $k$ vertices or an independent set on $l$ vertices.
(H. L. Abbott)
\end{theorem}

\begin{definition}[$C_{3} \lor C_{5}$ on the carrier \texttt{Fin 3 $\oplus$ Fin 5}]
\lean{SimpleGraph.grahamGraph}
\label{SimpleGraph.grahamGraph}
\leanok
\uses{SimpleGraph.join}
\texttt{C$_{3}$ $\lor$ C$_{5}$} on the carrier \texttt{Fin 3 $\oplus$ Fin 5}.
\end{definition}

\begin{definition}[The generalised Ramsey number $r(G_{1},\dots,G_m)$ over Mathlib's containment $\sqsubseteq$]
\lean{SimpleGraph.generalisedRamseyNumber}
\label{SimpleGraph.generalisedRamseyNumber}
\leanok
The generalised Ramsey number \texttt{r(G$_{1}$,\dots{},G\_m)} over Mathlib's
containment \texttt{$\sqsubseteq$}.

\textbf{7.2.6} Let $G_1, G_2, \ldots, G_m$ be simple graphs. The \emph{generalised
Ramsey number} $r(G_1, G_2, \ldots, G_m)$ is the smallest integer $n$ such that every
$m$-edge colouring $(E_1, E_2, \ldots, E_m)$ of $K_n$ contains, for some $i$, a subgraph
isomorphic to $G_i$ in colour $i$.
\end{definition}

\begin{theorem}[Trees embed in graphs of large minimum degree]
\lean{SimpleGraph.tree_isContained_of_minDegree_le}
\label{SimpleGraph.tree_isContained_of_minDegree_le}
Let $G$ be a finite simple graph and let $T$ be a tree on $m$ vertices. If every vertex
of $G$ has degree at least $m-1$, so that the minimum degree satisfies $\delta(G) \ge
m-1$, then $G$ contains a copy of $T$; that is, there is an injective map from the
vertices of $T$ into those of $G$ that sends every edge of $T$ to an edge of $G$,
exhibiting $T$ as a subgraph of $G$.
\end{theorem}

\begin{theorem}[Ramsey number of a tree versus a star]
\lean{SimpleGraph.generalisedRamsey_tree_star}
\label{SimpleGraph.generalisedRamsey_tree_star}
\uses{SimpleGraph.generalisedRamseyNumber}
Let $T$ be a tree on $m$ vertices, and let $n$ be a natural number for which $m-1$
divides $n-1$. Writing $K_{1,n}$ for the star on $n+1$ vertices — the complete bipartite
graph whose parts have sizes $1$ and $n$ — the two-colour generalised Ramsey number of
$T$ against $K_{1,n}$ satisfies
\[
r(T, K_{1,n}) = m + n - 1.
\]
\end{theorem}

\begin{theorem}[Chvátal's tree–complete graph Ramsey number]
\lean{SimpleGraph.chvatal_tree_complete}
\label{SimpleGraph.chvatal_tree_complete}
\uses{SimpleGraph.generalisedRamseyNumber}
Let $T$ be a tree on $m$ vertices, and let $n \ge 1$. Then the generalised Ramsey number
$r(T, K_n)$ — the least integer $N$ for which every colouring of the edges of the
complete graph $K_N$ with two colours contains either a copy of $T$ in the first colour
or a copy of the complete graph $K_n$ in the second — is given by
\[
r(T, K_n) = (m - 1)(n - 1) + 1.
\]
\end{theorem}

\begin{theorem}[Erdős's degree-majorisation of clique-free graphs]
\lean{SimpleGraph.degreeMajorised_of_cliqueFree}
\label{SimpleGraph.degreeMajorised_of_cliqueFree}
\uses{SimpleGraph.DegreeMajorised, SimpleGraph.degreeSequence}
(Erdős, 1970.) Let $m \ge 1$, and let $G$ be a finite simple graph that is
$K_{m+1}$-free, i.e.\ contains no complete subgraph on $m+1$ vertices. Then there exists
a complete multipartite graph $H$ on the same number of vertices such that $G$ is
degree-majorised by $H$; that is, writing the degree sequences of $G$ and $H$ in
nondecreasing order, each entry of $G$'s sequence is at most the corresponding entry of
$H$'s. Moreover, if $G$ and $H$ have the same degree sequence, then $G$ and $H$ are
isomorphic.
\end{theorem}

\begin{theorem}[Pointwise degree domination implies degree-majorisation]
\lean{SimpleGraph.degreeMajorised_of_forall_degree_le}
\label{SimpleGraph.degreeMajorised_of_forall_degree_le}
\uses{SimpleGraph.DegreeMajorised}
Let $G$ and $H$ be simple graphs on a common finite vertex set $V$. If every vertex has
at least as high a degree in $H$ as in $G$, that is $d_G(v) \le d_H(v)$ for all $v \in
V$, then $G$ is degree-majorised by $H$.
\end{theorem}

\begin{theorem}[Uniqueness in Turán's theorem]
\lean{SimpleGraph.turan_iso_of_card_edgeFinset_eq}
\label{SimpleGraph.turan_iso_of_card_edgeFinset_eq}
Let $V$ be a finite vertex set, let $m \geq 1$, and let $G$ be a simple graph on $V$
that contains no clique on $m+1$ vertices. Write $T$ for the Turán graph $T(|V|, m)$,
the complete $m$-partite graph on $|V|$ vertices whose parts differ in size by at most
one. If $G$ has exactly as many edges as $T$, then $G$ is isomorphic to $T$.
\end{theorem}

\begin{theorem}[Edge count forcing a triangle]
\lean{SimpleGraph.triangle_of_card_edgeFinset_gt}
\label{SimpleGraph.triangle_of_card_edgeFinset_gt}
Let $G$ be a finite simple graph on $n$ vertices with $m$ edges. If $n^2 < 4m$, then $G$
contains a triangle, that is, three mutually adjacent vertices.
\end{theorem}

\begin{theorem}[Non-bipartite dense graphs contain a triangle]
\lean{SimpleGraph.triangle_of_not_bipartite_of_card_edgeFinset_gt}
\label{SimpleGraph.triangle_of_not_bipartite_of_card_edgeFinset_gt}
Let $G$ be a finite simple graph that is not bipartite, with $n$ vertices and $m$ edges.
If \[(n-1)^2 + 4 < 4m,\] equivalently $m > \tfrac{(n-1)^2}{4} + 1$, then $G$ contains a
triangle, i.e. three pairwise adjacent vertices. (Erdős)
\end{theorem}

\begin{theorem}[Degree condition forcing a $K_{2,m}$ subgraph]
\lean{SimpleGraph.isContained_completeBipartite_two_of_sum_choose_gt}
\label{SimpleGraph.isContained_completeBipartite_two_of_sum_choose_gt}
Let $G$ be a finite simple graph with vertex set $V$, write $n = \abs{V}$ for its number
of vertices and $d(v)$ for the degree of a vertex $v$, and let $m \ge 2$ be an integer.
If
\[ (m-1)\binom{n}{2} < \sum_{v \in V} \binom{d(v)}{2}, \]
then $G$ contains a copy of the complete bipartite graph $K_{2,m}$.
\end{theorem}

\begin{theorem}[Edge bound forcing a $K_{2,m}$ subgraph]
\lean{SimpleGraph.isContained_completeBipartite_two_of_card_edgeFinset_gt}
\label{SimpleGraph.isContained_completeBipartite_two_of_card_edgeFinset_gt}
Let $G$ be a finite simple graph on a vertex set with $n$ vertices, and let $m \ge 2$ be
an integer. If the number of edges of $G$ satisfies
\[
\frac{(m-1)^{1/2}\, n^{3/2}}{2} + \frac{n}{4} < e(G),
\]
then $G$ contains the complete bipartite graph $K_{2,m}$ as a subgraph.
\end{theorem}

\begin{definition}[The unit-distance graph on a finite point set in the plane]
\lean{SimpleGraph.unitDistanceGraph}
\label{SimpleGraph.unitDistanceGraph}
The unit-distance graph on a finite point set in the plane.
\end{definition}

\begin{theorem}[Upper bound on unit-distance pairs in the plane]
\lean{SimpleGraph.card_unit_distance_pairs_le}
\label{SimpleGraph.card_unit_distance_pairs_le}
\uses{SimpleGraph.unitDistanceGraph}
Let $x_0, x_1, \dots, x_{n-1}$ be $n$ points of the Euclidean plane $\bbr^2$, and
consider the unit-distance graph on them, whose edges join the pairs of points lying at
distance exactly $1$. Then the number of such edges is at most
\[
\frac{n^{3/2}}{\sqrt{2}} + \frac{n}{4}.
\]
\end{theorem}

\begin{theorem}[The Kővári–Sós–Turán theorem]
\lean{SimpleGraph.isContained_completeBipartite_of_card_edgeFinset_gt}
\label{SimpleGraph.isContained_completeBipartite_of_card_edgeFinset_gt}
Let $G$ be a finite simple graph whose vertex set has $n$ elements, and write $e(G)$ for
its number of edges. Let $m \ge 1$ be an integer. If
\[
e(G) > \tfrac{1}{2}\,(m-1)^{1/m}\, n^{\,2 - 1/m} + \tfrac{1}{2}\,(m-1)\, n,
\]
then $G$ contains the complete bipartite graph $K_{m,m}$ as a subgraph.
\end{theorem}

\begin{theorem}[Schur's theorem on monochromatic solutions of x + y = z]
\lean{SimpleGraph.schur}
\label{SimpleGraph.schur}
\uses{SimpleGraph.ramseyTriangle}
Let $n \ge 1$, and let $r_n = r(3,\ldots,3)$ be the $n$-colour Ramsey number for
triangles. If $S_1, S_2, \ldots, S_n$ partition the set $\{1, 2, \ldots, r_n\}$ — that
is, each $S_i$ is a subset of $\{1, \ldots, r_n\}$ and every integer $k$ with $1 \le k
\le r_n$ lies in exactly one of the $S_i$ — then for some $i$ the part $S_i$ contains
integers $x, y, z$ (not necessarily distinct) with $x + y = z$.
\end{theorem}

\begin{definition}[The Schur number $s_{n}$]
\lean{SimpleGraph.schurNumber}
\label{SimpleGraph.schurNumber}
\leanok
The Schur number \texttt{s$_{n}$}. Nonemptiness of the \texttt{sInf} set is discharged
by Thm 7.10 (\texttt{s$_{n}$ $\le$ r$_{n}$}), itself deferred.

Let $s_n$ denote the least integer such that, in any partition of
$\{1,2,\ldots,s_n\}$ into $n$ subsets, there is a subset which contains a
solution to (7.19). [...] Also, from theorem 7.10 and exercise 7.2.3 we have the
upper bound
\[s_n\le r_n\le[n!\,e]+1\]
\end{definition}

\begin{theorem}[The third Schur number]
\lean{SimpleGraph.schurNumber_three}
\label{SimpleGraph.schurNumber_three}
\uses{SimpleGraph.schurNumber}
The Schur number for three parts equals fourteen, $s_3 = 14$. Equivalently, $14$ is the
least integer $N$ with the property that every partition of $\{1,2,\ldots,N\}$ into
three subsets has some subset containing elements $x,y,z$ (not necessarily distinct)
with $x+y=z$.
\end{theorem}

\begin{theorem}[A recursive lower bound for Schur numbers]
\lean{SimpleGraph.schurNumber_recursion}
\label{SimpleGraph.schurNumber_recursion}
\uses{SimpleGraph.schurNumber}
For each integer $n \ge 1$, the Schur numbers satisfy
\[
3\,s_{n-1} \le s_n + 1,
\]
where $s_m$ denotes the least $N$ such that every partition of $\{1, 2, \ldots, N\}$
into $m$ subsets has some part containing elements $x, y, z$ with $x + y = z$.
Equivalently, $s_n \ge 3 s_{n-1} - 1$.
\end{theorem}

\begin{theorem}[A lower bound for the Schur numbers]
\lean{SimpleGraph.schurNumber_lower_bound}
\label{SimpleGraph.schurNumber_lower_bound}
\uses{SimpleGraph.schurNumber}
For every integer $n \ge 3$, the Schur number $s_n$ satisfies
\[27\cdot 3^{\,n-3} + 1 \;\le\; 2\,s_n,\]
equivalently $s_n \ge \tfrac{1}{2}\bigl(27\cdot 3^{\,n-3} + 1\bigr)$.
\end{theorem}

\begin{definition}[The "far pairs" graph: adjacent iff distance exceeds $1/\sqrt 2$]
\lean{SimpleGraph.farGraph}
\label{SimpleGraph.farGraph}
\leanok
The "far pairs" graph: adjacent iff distance exceeds \texttt{1/$\sqrt{}$2}.

Let $G$ be the graph defined by
\[V(G) = \{x_1, x_2, \ldots, x_n\}\]
and
\[E(G) = \{x_i x_j \mid d(x_i, x_j) > 1/\sqrt{2}\}\]
where $d(x_i, x_j)$ here denotes the \emph{euclidean} distance between $x_i$ and
$x_j$.
\end{definition}

\begin{theorem}[Few pairs far apart in a diameter-one planar set]
\lean{SimpleGraph.card_farGraph_le}
\label{SimpleGraph.card_farGraph_le}
\uses{SimpleGraph.farGraph}
Let $x_1, x_2, \ldots, x_n$ be points in the Euclidean plane $\bbr^2$ whose set $\{x_1,
\ldots, x_n\}$ has diameter exactly $1$. Then the number of unordered pairs $\{x_i,
x_j\}$ whose Euclidean distance satisfies $d(x_i, x_j) > 1/\sqrt{2}$ is at most $\lfloor
n^2/3 \rfloor$.
\end{theorem}

\begin{theorem}[Diameter-one point set maximising far pairs]
\lean{SimpleGraph.exists_farGraph_card_eq}
\label{SimpleGraph.exists_farGraph_card_eq}
\uses{SimpleGraph.farGraph}
For every natural number $n$, there exist points $x_1, x_2, \ldots, x_n$ in the
Euclidean plane $\bbr^2$ such that the set $\{x_1, x_2, \ldots, x_n\}$ has diameter
exactly $1$ and exactly $\lfloor n^2/3\rfloor$ of the unordered pairs $\{x_i, x_j\}$
satisfy $d(x_i, x_j) > 1/\sqrt{2}$, where $d$ denotes Euclidean distance.
\end{theorem}

\begin{theorem}[Pannwitz's bound on unit-distance pairs in a set of diameter one]
\lean{SimpleGraph.card_unitDistanceGraph_le_of_diam_one}
\label{SimpleGraph.card_unitDistanceGraph_le_of_diam_one}
\uses{SimpleGraph.unitDistanceGraph}
Let $x_1, \dots, x_n$ be points of the Euclidean plane $\bbr^2$ whose set $\{x_1, \dots,
x_n\}$ has diameter exactly $1$. Then the number of unordered pairs $\{i, j\}$ with $i
\ne j$ and $\dist(x_i, x_j) = 1$ is at most $n$; equivalently, the unit-distance graph
on these points has at most $n$ edges.
\end{theorem}

\begin{definition}[The radio-range graph: adjacent iff within \texttt{range}]
\lean{SimpleGraph.radioGraph}
\label{SimpleGraph.radioGraph}
\leanok
The radio-range graph: adjacent iff within \texttt{range}.
\end{definition}

%%% added by the blueprint pipeline
\section{Additional declarations}

\begin{theorem}[Turán's edge bound for clique-free graphs]
\lean{SimpleGraph.turan_edge_bound}
\label{SimpleGraph.turan_edge_bound}
% not \leanok: the Lean proof body is currently sorry
Let $V$ be a finite vertex set, let $m$ be a natural number, and let $G$ be a simple
graph on $V$ that contains no clique on $m+1$ vertices. Then the number of edges of $G$
is at most the number of edges of the Turán graph $T(\abs{V}, m)$, the complete
$m$-partite graph on $\abs{V}$ vertices whose parts differ in size by at most one.
\end{theorem}
